\chapter{BASE TEÓRICA PARA A MONTAGEM DO LABORATÓRIO}

Este capítulo descreve as decisões metodológicas que estruturam o laboratório experimental. São justificados a escolha e as características do \textit{CICIDS2017} e do \textit{LycoS-IDS2017}, detalhado o \textit{pipeline} de preparação dos dados, definidas as três variantes de avaliação utilizadas nos experimentos e apresentadas as métricas adotadas para comparar os modelos.

O laboratório é organizado em dois níveis planejados desde o início. O primeiro é o experimento base (\textit{Baseline}), conduzido sobre o \textit{CICIDS2017} com os modelos \textit{Random Forest} e DNN nas três variantes de avaliação, e serve como referência de desempenho. O segundo nível é composto por oito cenários adicionais, concebidos para isolar o efeito de fatores específicos (qualidade do \textit{dataset}, redução de dimensionalidade e abordagem não supervisionada) sobre os resultados. Cada cenário varia exatamente um desses fatores em relação ao \textit{Baseline} ou ao cenário anterior, o que permite interpretar cada resultado como uma contribuição isolada e diretamente comparável.

\section{Escolha do \textit{dataset} e critérios de seleção}

A validade de um laboratório de detecção de intrusão depende diretamente da qualidade e da representatividade do conjunto de dados utilizado. Para que os resultados sejam comparáveis com a literatura e cubram uma variedade suficiente de comportamentos maliciosos, o \textit{dataset} deve conter tráfego de rede real, com rótulos confiáveis e diversidade de tipos de ataque \cite{neto_gomes}.

O \textit{CICIDS2017} (\textit{Canadian Institute for Cybersecurity Intrusion Detection System 2017}) foi escolhido como base para este laboratório. O conjunto foi gerado em ambiente controlado ao longo de cinco dias, com fluxos de rede capturados via \textit{CICFlowMeter}, uma ferramenta que extrai automaticamente atributos estatísticos de fluxos bidirecionais, como duração, volume de pacotes, tamanho médio de \textit{payload} e \textit{flags} TCP. Após consolidação dos arquivos diários, o conjunto totaliza 2.830.743 registros com 79 atributos por fluxo.

O \textit{dataset} cobre 15 classes: \textit{BENIGN} e 14 categorias de ataque, entre elas \textit{DDoS}, variantes de \textit{DoS} (\textit{Hulk}, \textit{GoldenEye}, \textit{slowloris}, \textit{Slowhttptest}), \textit{FTP-Patator}, \textit{SSH-Patator}, \textit{Heartbleed}, \textit{Infiltration}, \textit{PortScan} e três variantes de \textit{Web Attack} (\textit{Brute Force}, \textit{XSS} e \textit{Sql Injection}). Essa heterogeneidade intencional, com ataques volumétricos e raros representados no mesmo conjunto, é relevante para investigar a sensibilidade dos modelos ao desbalanceamento de classes, um problema central em IDS baseados em aprendizado de máquina \cite{cantone2024}.

A limitação mais documentada do \textit{CICIDS2017} são os erros de rótulo, as duplicatas e as inconsistências presentes nos arquivos originais, identificados por pesquisadores que analisaram criticamente o conjunto após sua publicação \cite{engelen2021}. Para isolar o efeito dessas inconsistências sobre o desempenho dos modelos, o laboratório inclui experimentos paralelos com o \textit{LycoS-IDS2017} \cite{engelen2021}, uma versão curada do mesmo conjunto obtida por correção de rótulos errôneos, normalização mais rigorosa e remoção de registros duplicados. O \textit{LycoS-IDS2017} contém 14 classes, sem a classe \textit{Infiltration}, e permite comparar os resultados em condições de dados mais confiáveis com o mesmo protocolo experimental, isolando a contribuição da qualidade do \textit{dataset} sobre o desempenho de cada modelo.

\section{\textit{Pipeline} de preparação de dados}

O \textit{CICIDS2017} é distribuído em arquivos CSV separados por dia de captura (segunda a sexta-feira). A primeira etapa do \textit{pipeline} foi a consolidação desses arquivos em um único conjunto de dados, seguida da remoção de registros com valores nulos ou infinitos, produzidos por divisões por zero durante a extração de atributos de fluxo pelo \textit{CICFlowMeter}. Após a limpeza, o conjunto resultou em 2.830.743 registros válidos.

A divisão treino-teste foi realizada de forma estratificada na proporção 70/30. A estratificação garante que a distribuição de classes no subconjunto de treino reproduza a distribuição original do \textit{dataset}, evitando que classes com baixo suporte fiquem ausentes ou sub-representadas na avaliação. Essa precaução é especialmente importante dado o severo desbalanceamento do \textit{CICIDS2017}, no qual a classe \textit{BENIGN} concentra mais de 80\% das amostras, enquanto classes como \textit{Heartbleed} e \textit{Web Attack -- Sql Injection} possuem apenas algumas dezenas de registros.

Para os experimentos com \textit{Deep Neural Network}, os atributos foram normalizados via \textit{MinMaxScaler}, que transforma cada atributo para o intervalo $[0, 1]$ dividindo cada valor pelo intervalo da coluna correspondente. Redes neurais são sensíveis à escala dos atributos: gradientes associados a atributos com maior amplitude tendem a dominar as atualizações de pesos durante a retropropagação, dificultando a convergência e o aprendizado de padrões em atributos de menor escala \cite{bochie2020sbrc}. O \textit{Random Forest}, por ser baseado em particionamentos recursivos por limiar, não é afetado pela escala dos valores e recebeu os dados na escala original.

Para os experimentos com redução de dimensionalidade (Cenários B, C, F e G), foi aplicada a técnica \textit{Mean Decrease in Impurity} (MDI). O MDI quantifica a importância de cada atributo pela sua contribuição média para a redução da impureza de Gini nas árvores de um \textit{Random Forest} previamente treinado no conjunto completo. Os atributos com menor contribuição média são descartados, e os modelos subsequentes são treinados e avaliados exclusivamente no espaço de \textit{features} reduzido, mantendo os mesmos conjuntos de treino e teste dos experimentos sem redução.

\section{Definição das variantes experimentais}

As três variantes foram construídas para separar dois objetivos distintos de avaliação: mensurar o desempenho dos modelos em ambiente fechado, onde todos os tipos de ataque estão representados no treinamento, e verificar a capacidade de generalização diante de ataques completamente ausentes no treinamento.

A Variante Completa inclui todas as 15 classes tanto no treinamento quanto no teste, resultando em 1.979.513 amostras para treino e 848.363 para teste. Ela serve como linha de base do desempenho máximo esperado de cada modelo e permite identificar fragilidades estruturais, em particular, a dificuldade de aprender classes com baixíssimo suporte mesmo quando todos os tipos de ataque estão representados.

A Variante Sem \textit{PortScan} remove a classe \textit{PortScan} do treinamento e a mantém exclusivamente no conjunto de teste, com 158.804 amostras (16,6\% do total de avaliação). O \textit{PortScan} é um ataque volumétrico de alta frequência: sua presença dominante no teste cria a variante mais severa de generalização, no qual os modelos precisam lidar com um tipo de ataque numericamente relevante para o qual não possuem nenhum exemplo de treinamento. O conjunto de treino nessa variante contém 1.868.545 amostras e o de teste 959.331.

A Variante Sem \textit{XSS} remove a classe \textit{Web Attack -- XSS} do treinamento, mantendo-a no teste com 652 amostras (0,08\% do total de avaliação). Diferentemente do \textit{PortScan}, o \textit{XSS} é raro e seus padrões de tráfego se sobrepõem com os de outros ataques Web, em especial o \textit{Brute Force}. Essa variante avalia a generalização em condições de baixíssimo suporte e alta ambiguidade de fronteira entre classes, permitindo observar como cada modelo distribui amostras desconhecidas entre as classes mais próximas.

A combinação das duas variantes de generalização permite caracterizar o comportamento dos modelos em dois extremos complementares: um ataque ausente no treino com alta presença estatística no teste (\textit{PortScan}) e um ataque ausente no treino que é raro e semanticamente ambíguo (\textit{XSS}).

\section{Configuração dos modelos e ambiente computacional}

Os três modelos foram implementados em Python em ambiente local, utilizando \textit{Scikit-learn} para o \textit{Random Forest} e \textit{TensorFlow}/\textit{Keras} para a DNN e o \textit{Autoencoder}.

O \textit{Random Forest} foi configurado com 100 árvores de decisão (\texttt{n\_estimators=100}), semente aleatória fixada em 42 (\texttt{random\_state=42}) e paralelização total dos núcleos disponíveis (\texttt{n\_jobs=-1}). Os demais parâmetros seguiram os valores padrão da biblioteca.

A DNN foi construída com duas camadas ocultas: a primeira com 128 neurônios e a segunda com 64, ambas com ativação ReLU e \textit{dropout} de 0,2. A camada de saída possui um neurônio por classe com ativação \textit{softmax}. O modelo foi compilado com otimizador \textit{Adam} e função de perda \textit{sparse categorical crossentropy}, treinado por 10 épocas com \textit{batch size} de 256 amostras e 10\% do conjunto de treino reservado para validação interna.

O \textit{Autoencoder} adota arquitetura simétrica: o codificador reduz a entrada de 78 atributos para 64, 32 e 16 dimensões; o decodificador reconstrói o caminho inverso (16 $\to$ 32 $\to$ 64 $\to$ 78). As camadas ocultas utilizam ativação ReLU e a camada de saída utiliza sigmoide. O modelo foi compilado com otimizador \textit{Adam} e perda MSE, treinado por até 50 épocas com \textit{batch size} de 256, fração de validação de 10\% e \textit{EarlyStopping} com paciência de 3 épocas, restaurando os melhores pesos ao término.

As especificações completas do ambiente computacional utilizado nos experimentos são apresentadas na Tabela~\ref{tab:ambiente}.

\begin{table}[H]
\centering
\caption{Especificações do ambiente computacional}
\label{tab:ambiente}
\begin{tabular}{ll}
\hline
Componente & Especificação \\
\hline
CPU             & AMD Ryzen 5 5600 6-Core (3,50 GHz) \\
GPU             & NVIDIA GeForce RTX 3060 (12 GB) \\
RAM             & 32 GB \\
Sistema Operacional & Microsoft Windows (64-bit) \\
Python          & 3.12.10 \\
TensorFlow      & 2.21.0 \\
Keras           & 3.13.2 \\
Scikit-learn    & 1.8.0 \\
Pandas          & 3.0.1 \\
NumPy           & 2.4.3 \\
\hline
\end{tabular}
\end{table}

\section{Métricas e critérios de avaliação}

A avaliação dos modelos utilizou métricas por classe e métricas agregadas, complementadas pelo tempo de inferência como estimativa do custo computacional em produção.

Precisão mede a fração de predições de uma classe que correspondem efetivamente a amostras dessa classe: $P = VP\,/\,(VP + FP)$, onde $VP$ são os verdadeiros positivos e $FP$ os falsos positivos. Uma precisão baixa em uma classe de ataque indica que o modelo está gerando muitos falsos alarmes para aquele tipo de ameaça.

Revocação mede a fração de amostras de uma classe que foram corretamente identificadas pelo modelo: $R = VP\,/\,(VP + FN)$, onde $FN$ são os falsos negativos. Para um IDS, revocação baixa em uma classe de ataque significa que o modelo está deixando passar ataques reais, o erro mais crítico do ponto de vista operacional \cite{santos2023ercemapi}.

O \textit{F1-score} é a média harmônica entre precisão e revocação: $F1 = 2 \cdot P \cdot R\,/\,(P + R)$. A média harmônica penaliza desequilíbrios acentuados entre as duas métricas, de modo que um F1 elevado exige desempenho satisfatório em ambas.

A média macro calcula a média aritmética do F1 de todas as classes, atribuindo peso igual a classes raras e majoritárias. Em \textit{datasets} com desbalanceamento severo como o \textit{CICIDS2017}, a média macro é a principal métrica de comparação entre modelos, pois expõe dificuldades em classes de baixo suporte que as métricas ponderadas tendem a mascarar.

A média ponderada calcula a média do F1 ponderada pelo número de amostras de cada classe. Ela reflete o desempenho sobre o volume total do tráfego e é dominada pela classe \textit{BENIGN} dada sua representatividade, sendo interpretada como complemento, e não substituto, da média macro.

A acurácia global é a fração total de amostras classificadas corretamente. Em cenários desbalanceados, pode permanecer elevada mesmo quando o modelo falha sistematicamente em classes minoritárias, razão pela qual é sempre analisada em conjunto com as métricas por classe.

A matriz de confusão detalha, para cada classe real, a distribuição das predições entre todas as classes possíveis. Ela permite identificar padrões específicos de erro, como ataques de uma classe sendo absorvidos pela classe \textit{BENIGN} ou classes confundidas entre si, que métricas escalares não capturam.

O tempo de inferência foi medido sobre o conjunto de teste integral e adotado como estimativa do custo computacional em ambiente de produção. Modelos com menor tempo de inferência são mais adequados para redes com alto volume de tráfego onde a latência de detecção é um requisito operacional \cite{sbrc_minicursos}. Com o protocolo experimental completamente definido, abrangendo dados, \textit{pipeline}, cenários e métricas, o capítulo seguinte apresenta e analisa os resultados obtidos em cada um dos dez experimentos realizados.
